\section{Supporting evidence on one canonical circle}\label{sec:evidence}
Take \(Y=G_{24}\), \(Z=F\), \(\Gamma=\{|z-36|=9\}\), and \(0\le t\le T=10^{-4}\). The measured quantity is
\(\Delta(t)=D_F(t)-D_{24}(t)\). The earlier complete-theta certificate \cite{hdc} counts two zeros in the rectangle \((26,46)+i(-2,2)\) at time zero. Since the time-zero xi zeros lie in \(|\Im z|\le1\), every possible zero in either disk of radius \(9\pm1/64\) is in that rectangle. Two real sign changes exhaust its count. The following additional bounds are needed to extend this geometry to the slab and determine the sign of the correction \cite{jh07,jh08}.

The construction in Section~\ref{sec:time} used \(\rho=1/20\), \(R_0=24\), \(R=10\) and time degree \(K=10\). Initial moments through order 66 supply the moving head through that order. Its Gram perturbation is below \(0.051272\). With 144 centered spatial terms for the source and 192 origin exponential terms for the head, the respective joint tails on radius ten are below \(4.750\times10^{-29}\) and \(8.876\times10^{-27}\). Outward coefficient intervals include all rounding. The polynomial for \(F-G_{24}\) is formed before evaluation over the time interval.

On both collar boundaries a Taylor enclosure covers each of 512 angular sectors. Nonvanishing for every \(t\) preserves the source counts by homotopy, and the strict relative bound transfers both counts to the head. Sign changes in \((28,29)\) and \((42,43)\) for each field at every \(t\) place both simple real roots inside radius eight. Equality of the two circle counts excludes zeros from the annulus. Maximum modulus applied to the quotient and reciprocals therefore gives the uniform bounds
\[
\begin{gathered}
  |F|>1.4331937907\times10^{-5},\qquad
  |G_{24}|>1.4331935452\times10^{-5},\\
  |(F-G_{24})/G_{24}|<1.715356\times10^{-7}.
\end{gathered}
\]
In particular \(b=1\). The exact choice \(q_0=q\) gives zero modeling error \(e\). Its numerical evaluation errors are included in the outward intervals. The two terms of the nonlinear allowance in Theorem~\ref{th:signed} sum to less than \(2.846816\times10^{-12}\).

\begin{table}[htb]
\centering\small
\caption{Certified enclosures of the same correction over the entire interval. Widths use the stored endpoints before outward decimal rounding.}\label{tab:certificate}
\begin{tabular}{@{}lcc@{}}
\toprule
Method & Interval for \(10^8\Delta(t)\) & Bound on width of \(\Delta\)\\
\midrule
Signed contour & \((-4.958280826,-4.956734690)\) & \(1.546135\times10^{-11}\)\\
Direct displacement & \((-4.957996104,-4.957019318)\) & \(9.767860\times10^{-12}\)\\
\bottomrule
\end{tabular}
\end{table}

Table~\ref{tab:certificate} uses \eqref{eq:residue} and \eqref{eq:direct} with the same source, head, root intervals and numerator. In particular, the direct test retains the correlation between the two fields at each time. The slope bounds on both connecting segments exclude zero. The unsigned logarithm bound allows \(|\Delta|<2.624495\times10^{-5}\), leaving the sign unresolved, although it already proves positivity of the absolute gap on this known chart. The signed and direct evaluations show that the head overestimates the squared half-gap throughout the interval. The direct enclosure lies strictly inside the signed one.

The time-zero calculation alone encloses \(\Delta(0)\) near \(-4.957710484497283\times10^{-8}\) with direct width below \(2.073\times10^{-35}\). Subtracting this saved interval from the same-time direct evaluation at \(T\) gives
\[
 4.05561835\times10^{-12}<\Delta(T)-\Delta(0)<4.05588776\times10^{-12}.
\]
This measures a decrease in the magnitude of the error between the two endpoints. Its derivative in the interior has no certified sign here. The full-interval enclosure came from the range of a joint analytic model. The repository records the outward endpoints, evaluator versions, counts, and the matched internal replay \cite{jh07,jh08}.

A simple event check distinguishes the two possible sign decisions. Take
\(Y=(z-36)^2-2t-1\) and \(Z=(z-36)^2-2t\). Both obey the unnormalized heat equation, and their squared half-gaps are \(2t+1\) and \(2t\). The correction is exactly \(-1\), even at the target double at \(t=0\). The strict absolute-gap condition following Theorem~\ref{th:signed} fails there, as do distinct target brackets and a target slope floor. The contour coordinates remain defined. This example isolates the distinction between an approximation error and a collision margin.
